\section{Visit coefficients in scan order}
The quantizer's level array uses natural row-major order: index $8v+u$ identifies
vertical frequency $v$ and horizontal frequency $u$. Coefficient coding visits
these positions in a zigzag order that tends to encounter lower frequencies first.
The order is a permutation of sixty-four positions. It changes the sequence in
which values are described; it does not apply another transform or move samples
between pictures.

\begin{center}
  \tagpdfsetup{table/header-rows={1}}
  \begin{tabular}{rrrrr}
    \textbf{Scan position} & \textbf{$v$} & \textbf{$u$} & \textbf{Natural index} & \textbf{Byte offset} \\
    0                      & 0            & 0            & 0                      & 0                    \\
    1                      & 0            & 1            & 1                      & 2                    \\
    2                      & 1            & 0            & 8                      & 16                   \\
    3                      & 2            & 0            & 16                     & 32                   \\
    4                      & 1            & 1            & 9                      & 18                   \\
    5                      & 0            & 2            & 2                      & 4
  \end{tabular}
\end{center}
The first two AC positions therefore occupy natural indices one and eight.
Iterating ordinary row indices instead would visit natural index two before
eight and describe a different coefficient block. The supplied table contains
all sixty-four natural indices in scan order. Each index addresses a halfword,
so multiply it by two before adding the level-array base.

For INTRA, write the separate eight-bit DC code first, then begin ordinary events
at scan position one. For INTER, begin at scan position zero; DC is an ordinary
level there. These rules match the quantizer's earlier distinction between a
special INTRA DC code and ordinary signed levels.

\section{Describe zeros with RUN and finish with LAST}
A coefficient event describes a nonzero ordinary level and the number of zeros
immediately before it in scan order. This zero count is RUN. Its origin is the
previous emitted nonzero, or the beginning of the ordinary scan if there has
not been an event yet. A nonzero resets the run count. Zeros after the final
nonzero are implicit and do not require another event.

LAST is one on the final nonzero event and zero on earlier events. It means
that no further nonzero coefficients follow, not that this event occupies scan
position sixty-three. Find the final nonzero scan position before emitting, or
retain enough look-ahead to determine that fact correctly. If every ordinary
level is zero, emit no ordinary events at all. The macroblock layer will later
signal whether a block has a TCOEF payload; there is no invented end-of-block
marker for an empty payload here.

For an INTRA block with DC code 255, level one at natural index one and level
minus one at natural index eight, the ordinary scan begins with those two
nonzeros. Both RUN values are zero. The first event has LAST zero; the second
has LAST one. The remaining sixty-one AC positions are zero and are implicit.

\begin{center}
  \tagpdfsetup{table/header-rows={1}}
  \begin{tabular}{rrrrl}
    \textbf{Order} & \textbf{LAST} & \textbf{RUN} & \textbf{LEVEL} & \textbf{Meaning}                        \\
    First          & 0             & 0            & 1              & More nonzeros follow                    \\
    Second         & 1             & 0            & $-1$           & Final nonzero; remaining positions zero
  \end{tabular}
\end{center}
With only a final INTER nonzero at scan position sixty-three, RUN is sixty-three.
For INTRA, DC is excluded, so the corresponding final AC event has RUN sixty-two.
Both fit the six-bit escape RUN field. Counting a nonzero itself as part of RUN
would shift every subsequent decoded coefficient.

\section{Choose a compact code or a baseline escape}
The baseline table assigns compact variable-length codes to 102 event combinations
of LAST, RUN and magnitude. The final code bit carries sign: zero for positive,
one for negative. The supplied numeric code table stores that sign bit as zero;
OR one into it for a negative level. Its width includes the sign bit.

\begin{center}
  \tagpdfsetup{table/header-rows={1}}
  \begin{tabular}{rrrrl}
    \textbf{LAST} & \textbf{RUN} & \textbf{LEVEL} & \textbf{Bits} & \textbf{Code}  \\
    0             & 0            & 1              & 3             & \texttt{100}   \\
    0             & 0            & $-1$           & 3             & \texttt{101}   \\
    1             & 0            & 1              & 5             & \texttt{01110} \\
    1             & 0            & $-1$           & 5             & \texttt{01111}
  \end{tabular}
\end{center}
Every other legal ordinary event can use the twenty-two-bit escape form:
\[
  \underbrace{\texttt{0000011}}_{7\ \mathrm{bits}}\quad
  \underbrace{\mathrm{LAST}}_{1}\quad
  \underbrace{\mathrm{RUN}}_{6}\quad
  \underbrace{\mathrm{LEVEL}}_{8}.
\]
The eight-bit LEVEL is the signed two's-complement value, so minus 127 is
\texttt{10000001}. Zero and minus 128 are forbidden in the baseline ordinary
escape field. This differs from a compact code's separate sign bit. Convert the
signed level to its low eight bits only after validating its legal nonzero range.

The escape form is also legal for combinations that have compact codes. Thus
correct encoders can produce different bits for the same levels. The compact
reference emits the worked block as eight DC bits, three first-event bits and
five last-event bits: sixteen bits total. An escape-only alternative emits the
same DC followed by two twenty-two-bit events: fifty-two bits. Both decode to
the same level array. The complete encoder may choose either valid syntax;
this exercise does not define correctness as matching a private byte string.

The numeric tables are drawn from H.263 Table 16 and its scan figure. Their
provenance is retained with the data. The source includes the complete baseline
table, not only the four worked entries above. A linear lookup is sufficient for
initial construction; its search work and emitted-bit savings can be measured
before replacing it with a direct lookup structure.

\section{Plan a complete block before touching caller output}
A complete coefficient block can require many append calls. Field-level preflight
inside the bit writer cannot make that sequence atomic. Build a private plan of
code values and widths first, validating all levels as you go. Sum the widths of
the chosen representation, then compare that complete total with the caller's
remaining bit capacity. If it does not fit, return before reading or writing the
caller buffer. Different valid choices may have different space requirements.

For $N$ ordinary nonzeros, every legal representation in this stage needs at most
$22N$ bits, plus eight for INTRA DC. When that bound fits, the function must be
able to succeed. With less space, a compact plan might still fit while an
escape-only plan does not. A capacity failure describes the selected plan;
it does not prove that no other legal representation could fit.

Retain the destination byte window privately while emitting the plan. Copy its
initial bytes so existing prefix and unused suffix bits survive the append.
Give the bit writer a private context starting at the original bit offset within
its first byte. If any child fails, abandon the private work and retain both
caller bytes and caller count. Only after every child succeeds should you copy
the completed touched bytes back and commit the caller's new used count.

The largest INTER payload has 64 escape events, or 1408 bits. Starting at bit
offset seven, it touches at most $\lceil(1408+7)/8\rceil=177$ bytes. INTRA has at
most 63 ordinary events plus eight DC bits and touches at most 176 bytes. A
192-byte private window therefore suffices for either, including an unaligned
start. A plan of at most 64 value/width records covers both cases as well.

\section{Measure code selection separately from bit packing}
Use the same byte-fragment child when comparing compact lookup with escape-only
coding. Otherwise a faster bit writer could be mistaken for a cheaper code-selection
algorithm. The named worked fixture is the earlier DC 255 with two AC events.
The dense fixture has DC 255 and sixty-three ordinary levels alternating minus
one and plus one in natural order. Both start at bit offset three.

\begin{center}
  \tagpdfsetup{table/header-rows={1}}
  \begin{tabular}{llrr}
    \textbf{Fixture} & \textbf{Coding} & \textbf{Emitted bits} & \textbf{Guest instructions} \\
    Worked           & Compact lookup  & 16                    & 3206                        \\
    Worked           & Escape only     & 52                    & 2952                        \\
    Dense            & Compact lookup  & 199                   & 12861                       \\
    Dense            & Escape only     & 1394                  & 18550
  \end{tabular}
\end{center}
These complete entry-to-return observations include validation, scan discovery,
plan construction, private staging, child appends and publication. Both native
byte orders give the same counts. The worked compact form saves bits while its
linear lookup adds work. On the dense fixture, shorter appends outweigh that
lookup cost and reduce both measures. A direct index for LAST, RUN and magnitude
is a possible next lookup experiment; it must retain the same decoded meaning
and be measured with the same child writer and fixture.

The measurements are guest instruction work, not host wall time or complete-encoder
throughput. The full encoder's clip size and optimization priorities still need
the completed picture path, file output and independent decoding measurements.
